% TOP_PROBLEM: {"id":30007614,"problem_number":"LOCAL-30007614","title":"Uncovered-interest-parity deviations","category_id":21,"category":{"id":21,"name":"economics","display_name":"Economics","description":"Open economic theory statements, published research agendas, and proposed scoped specifications; question provenance and historical priority are qualified per record.","slug":"economics","order_index":21,"created_at":"2026-10-08T00:00:00Z"},"status":"research_question","external_url":null,"legacy_ids":[],"published":false,"statement":"In the explicitly specified setting: How much of currency excess returns reflects risk compensation, belief errors or segmented financial markets? The mathematical statement above fixes the target, assumptions and scope of this catalogue formulation.","economics":{"original_id":103,"field":"International finance and exchange rates","kind":"identification","formalization_basis":"adapted_research_question","source_status":"explicit_open","priority_status":"earliest_found","priority_note":"Fama cites earlier forward-rate and currency-premium research, including Hansen–Hodrick1980 and Bilson1981. It is a classic verified decomposition, not proven first discovery.","earliest_verified_reference":{"authors":["Eugene F. Fama"],"title":"Forward and Spot Exchange Rates","year":1984,"url":"https://eml.berkeley.edu/~craine/EconH195/Fall_14/webpage/Fama_Forward%20Discount.pdf","doi":"10.1016/0304-3932(84)90046-1","locator":"Introduction, p.319; Section 2, pp.320–323; Section 5, pp.334–336","evidence_summary":"Compares incomplete, nonexclusive explanations of currency-premium behavior."},"current_reference":{"authors":["Ece Ozge Emeksiz","Andres Fernandez","Nikhil Patel","Ivan Petrella","Tatjana Schulze"],"title":"Drivers of Exchange Rates in EMDEs: Implications for Foreign Exchange Intervention","year":2026,"url":"https://www.imf.org/en/publications/staff-discussion-notes/issues/2026/09/16/drivers-of-exchange-rates-in-emdes-implications-for-foreign-exchange-intervention-578715","doi":"10.5089/9798229055208.006","locator":"Official overview: Conceptual framework and methodology; Stylized facts and empirical findings","evidence_summary":"Context-specific fundamental/financial shock decomposition; full PDF not audited."},"formalization_note":"The parity condition and decomposition are standard, and many subcases are studied. Observable currency returns alone do not identify risk, beliefs or segmentation contributions."},"statusQualification":"Published question or research agenda verified at the cited locator. The mathematical specification is adapted for this catalogue and includes additional assumptions; source publication does not establish first-ever priority or certify this exact model remains unsolved. The parity condition and decomposition are standard, and many subcases are studied. Observable currency returns alone do not identify risk, beliefs or segmentation contributions.","statusReviewedAt":"2026-10-08"}
% FIRST_POSTED: 2026-10-08
% ENTRY_KIND: statement_only
% !TeX program = lualatex
\documentclass[11pt]{article}
\usepackage{fontspec}
\usepackage[a4paper,margin=25mm]{geometry}
\usepackage{amsmath,amssymb,mathtools}
\usepackage{xurl}
\usepackage[hidelinks]{hyperref}
\newcommand{\sref}[2]{\hyperlink{source:#1:#2}{[S#2]}}
\setlength{\emergencystretch}{3em}
\begin{document}
% problemId:30007614
\section*{Uncovered-interest-parity deviations}\label{30007614}
\subsection{Definitions and mathematical statement}
Let \(s_t\) be log domestic-currency units per foreign currency, and \(i_t,i_t^*\) comparable one-period domestic and foreign interest rates. Define the first-order log excess return from borrowing domestically and investing abroad as \(x_{t+1}=i_t^*-i_t+s_{t+1}-s_t\). Let \(\mathbb P\) be the objective probability law and \(\mathcal F_t\) dated available information. Under this conditional law, its expectation is the currency premium \(\pi_t=E_{\mathbb P}[x_{t+1}\mid\mathcal F_t]\); uncovered interest parity requires \(\pi_t=0\). An investor's subjective law \(\mathbb Q_t\) gives an expectation wedge
\[b_t=E_{\mathbb Q_t}[s_{t+1}\mid\mathcal F_t]-E_{\mathbb P}[s_{t+1}\mid\mathcal F_t].\]
The task is to identify which part of predictable excess returns is attributable to compensation for priced risks, belief errors, or trading and financing constraints. Specify a discount-factor model, a survey measurement model for \(\mathbb Q_t\), and the mechanism relating intermediary constraints to attainable portfolios. Determine the identifying restrictions separating these components and derive bounds when objective expectations or marginal investors are unobserved. Evaluate variation by currency, horizon and financial state. Realized excess returns may remain predictable under rational expectations with time-varying risk compensation; a regression slope inconsistent with parity does not by itself identify irrationality. Exact gross-return analysis should be used when the log approximation materially changes conclusions.

\par\textbf{Formulation and scope.} The parity condition and decomposition are standard, and many subcases are studied. Observable currency returns alone do not identify risk, beliefs or segmentation contributions.

\par\textbf{Open-question provenance.} An explicit question or research agenda was verified at the source locator. This does not by itself certify that every later version remains unresolved \sref{30007614}{1}.

\par\textbf{Historical priority.} Fama cites earlier forward-rate and currency-premium research, including Hansen–Hodrick1980 and Bilson1981. It is a classic verified decomposition, not proven first discovery.

\subsection{Short English statement}
In the explicitly specified setting: How much of currency excess returns reflects risk compensation, belief errors or segmented financial markets? The mathematical statement above fixes the target, assumptions and scope of this catalogue formulation.

\subsection{Sources}
\begin{itemize}\raggedright
\item[\textnormal{[S1]}]\hypertarget{source:30007614:1}{} Eugene F. Fama. 1984. Forward and Spot Exchange Rates. Introduction, p.319; Section 2, pp.320–323; Section 5, pp.334–336.
\url{https://eml.berkeley.edu/~craine/EconH195/Fall_14/webpage/Fama_Forward%20Discount.pdf}.
DOI: 10.1016/0304-3932(84)90046-1.
{\small\textit{Earliest verified question/agenda reference; historical priority qualified below. Compares incomplete, nonexclusive explanations of currency-premium behavior.}}

\item[\textnormal{[S2]}]\hypertarget{source:30007614:2}{} Ece Ozge Emeksiz, Andres Fernandez, Nikhil Patel, Ivan Petrella, Tatjana Schulze. 2026. Drivers of Exchange Rates in EMDEs: Implications for Foreign Exchange Intervention. Official overview: Conceptual framework and methodology; Stylized facts and empirical findings.
\url{https://www.imf.org/en/publications/staff-discussion-notes/issues/2026/09/16/drivers-of-exchange-rates-in-emdes-implications-for-foreign-exchange-intervention-578715}.
DOI: 10.5089/9798229055208.006.
{\small\textit{Supporting or current literature; see evidence qualification. Context-specific fundamental/financial shock decomposition; full PDF not audited.}}
\end{itemize}
\end{document}
